\documentclass[11pt]{article}
\usepackage{amsmath,amssymb}
\title{An actual nonempty exchange compensation with a summable bound}\author{}\date{}
\begin{document}\maketitle
\paragraph{Input and target.}
The raw temporal-tree obstruction of the raw tree quartic energy obstruction calculation is used,
including its six exact plaquette checks. It concerns raw coordinate
vertices, not the completed physical coefficient. Here a SPECIFIED
nonempty conditional integration yields a genuinely decreasing bound
for its complete classical quartic plus exchange. This is not a bound
on the quantum coefficient $C_R''$ or the original nonlinear remainder.

\paragraph{1. Original record and complete conditional fiber.}
Take the same side-$R$ cube, $R\ge6$, with identity outer boundary.
Retain ALL temporal links and integrate ALL free spatial links with
original Haar measure and every incident Wilson plaquette. On one
interior spatial square place
\[
 \phi_0(2,2,2)=e_1,\quad \phi_0(3,2,2)=e_2,\quad
 \phi_0(3,3,2)=-e_1,\quad \phi_0(2,3,2)=-e_2,
\]
and put $\phi_0=0$ elsewhere. Let $\phi=R^{-1/2}\phi_0$.
The distributed temporal record is
$U_0(t,x)=\exp(\alpha\phi(x))$ for every temporal layer; the free
spatial classical minimum has first tangent $a_i=0$.
Indeed the temporal plaquette curls are independent of time, so the
linear spatial force cancels between adjacent temporal plaquettes.
The original energy is $E(a)=24$.

The original interior gauge
$g(t,x)=\exp(\alpha t\phi(x))$ for $0<t<R$, with $g=1$ on the
outer boundary, maps this temporal record to
\[
 U_0^{\rm ax}(t,x)=1\quad(t<R-1),\qquad
 U_0^{\rm ax}(R-1,x)=\exp(\alpha R\phi(x)).                    \tag{1}
\]
It also maps each integrated spatial link by its honest endpoint
gauges. Both conditional fibers and their entire original action
are mapped bijectively, with Haar Jacobian one. Their RAW conditional
partition functions are equal, without dropping a determinant or
normalizer. If a local chart is used its cutoff must be transported;
no independent fixed canonical cutoff is inserted.

Consequently their reduced CLASSICAL minimum actions agree for each
small $\alpha$ at fixed $R$. The axial first tangent is precisely the
temporal-tree type $h_i=-t\delta_i\phi$, $h_0=0$ except in the last
layer. The second-order minimum displacement is not discarded.

\paragraph{2. Literal cubic force and physical quartic.}
Use the normalization $s=1-\tfrac12\operatorname{Tr}U_p$ so that the
quadratic action is $E/2$. In the distributed record only temporal
links have nonzero first tangent. The two-link plaquette product gives
\[
 T_4(a)=-2/R.                                                  \tag{2}
\]
There are 16 spatial edges with one nonzero endpoint and four with two
orthogonal unit endpoints. Before the common $R/R^2$ scaling, their
quartics are respectively $-1/24$ and $-1/3$, totaling $-2$.

The ORIGINAL cubic force on a free spatial link is
\[
 j_i(t,x)=\nabla_{I,i,t,x}T_3(a)
       =2\phi(x)\mathbin\times\phi(x+\hat i)=R^{-1}j_{0,i}(x).
                                                               \tag{3}
\]
Each of its two adjacent temporal plaquettes contributes once. On the
chosen four-cycle, $j_0$ is exactly twice the oriented boundary of that
single spatial plaquette, in color three:
\[
 j_0=\pm2d_3^* e_p,\qquad D_3^*j_0=0.                         \tag{4}
\]
Signs depend only on orientation and do not affect the estimate.
There is no longitudinal forcing hidden in (4).

\paragraph{3. Exact exchange and its uniform bound.}
With temporal links fixed, the original spatial quadratic precision is
\[
 A_I=\Delta_t^D\otimes I+I\otimes H_3,
 \qquad H_3=d_3^*d_3.                                        \tag{5}
\]
Spatial one-forms have their actual relative boundary condition;
time has Dirichlet endpoints because the temporal boundary spatial
links are fixed. $A_I$ is strictly positive, even on spatial gradient
modes. On the co-closed source (4), the variational comparison gives
\[
 \begin{split}
 0\le j^*A_I^{-1}j
 &\le\frac{R-1}{R^2}\langle j_0,H_3^+j_0\rangle\\
 &=\frac{4(R-1)}{R^2}
       \langle e_p,d_3H_3^+d_3^*e_p\rangle
 \le\frac{4(R-1)}{R^2}\le\frac4R .
 \end{split} \tag{6}
\]
The last operator is the exact orthogonal projection onto the range
of the original spatial curl. No infinite-volume Green function or
unproved spectral gap is used. Equivalently one may use the variational
supremum for $j^*A_I^{-1}j$, discard the nonnegative temporal energy,
and maximize in each slice; $j_0$ annihilates the spatial kernel.

The conditional classical minimum therefore has the exact quartic
coefficient
\[
 T_{4,\rm red}=T_4(a)-\tfrac12j^*A_I^{-1}j,
 \qquad \boxed{-4/R\le T_{4,\rm red}\le-2/R.}                 \tag{7}
\]
By the actual gauge equivalence in paragraph1 this is ALSO
\[
 T_4(h)-\tfrac12\nabla_I T_3(h)^*A_I^{-1}\nabla_I T_3(h)         \tag{8}
\]
in the concentrated temporal-tree record. On the chosen spatial square the axial plaquette has zero linear curl
but quartic $8t^4/R^2$; its sum is of order $R^3$, while the last
temporal layer costs at most $CR^2$. Thus the coordinate growth
of the raw axial vertex is compensated in this genuine nonempty
conditional expression. This is not an assertion that the separate
exchange or raw vertex is small in that frame.

For dyadic block sizes $R_j=2^jR_0$, the absolute bounds in (7) sum to
at most $8/R_0$. In canonical action units this quartic contribution
is multiplied by $B^{-1}$. This is a one-record scale bound, NOT a
claim that a physical-volume cover can be summed without block counts.
The chosen retained temporal record and all exterior data have been
specified; no probability of this record is asserted.

\paragraph{4. Verification and remaining main estimate.}
The finite calculation checks the energy $24$, $-2$ unscaled physical quartic,
all four cubic forces, and their zero divergence. A separate original
spatial incidence solve plus exact temporal sine decomposition gives
$R T_{4,\rm red}\simeq-3.02579,-3.10034,-3.17475$ at $R=6,8,12$,
inside the PROVED interval $[-4,-2]$. Those numerical solves are only
checks; (6) proves the uniform estimate.

The equality of conditional partition functions retains all Haar,
quantum determinant and vacuum scalar factors, but (7) estimates ONLY
the reduced classical quartic/exchange contribution. It does not bound
those other factors, the full induced cubic variance under the joint
skeleton law, or the physical first nonlinear background response
$C_R''$. For general records the cubic force can have a spatial
longitudinal part, so the projection argument (6) cannot be assumed.
This example identifies an actual cancellation and a sufficient curl
structure, not a general Sobolev replacement or S94 closure.
\end{document}
